\chapter{Технологический раздел}

\section{Используемое программное обеспечение}
Для реализации проекта была выбрана реляционная СУБД PostgreSQL.

Серверная часть CLI-приложения реализована на C\# (.NET 8) с использованием ASP.NET Core, EF Core и гексагональной архитектуры для обеспечения модульности, тестируемости и чёткого разделения ответственностей.

\section{Реализация сущностей системы}
В листингах~\ref{lst:tbl1}--\ref{lst:tbl3} приведены запросы для создания таблиц.

\begin{lstlisting}[label=lst:tbl1,caption=Запросы для создания таблиц (часть 1)]
CREATE TABLE IF NOT EXISTS client (
	id              SERIAL PRIMARY KEY,
	client_login    login_inst,
	client_password password_inst,
	mail            mail_inst,
	privilege       role_type
);

CREATE TABLE IF NOT EXISTS chart (
	id              SERIAL PRIMARY KEY,
	city            VARCHAR(255),
	title           VARCHAR(255),
	svg_content     XML
);

CREATE TABLE IF NOT EXISTS branch (
	id              SERIAL PRIMARY KEY,
	title           VARCHAR(255),
	color           decimal_hexcolor,
	access          access_type
);

CREATE TABLE IF NOT EXISTS station (
	id              SERIAL PRIMARY KEY,
	duty_id         INT,
	title           VARCHAR(255),
	occupancy       occupancy_level,
	access          access_type,
	open_time       TIME,
	close_time      TIME
);
\end{lstlisting}

\clearpage

\begin{lstlisting}[label=lst:tbl2,caption=Запросы для создания таблиц (часть 2)]
CREATE TABLE IF NOT EXISTS transition (
	id              SERIAL PRIMARY KEY,
	duty_id         INT,
	occupancy       occupancy_level,
	access          access_type,
	duration        INTERVAL,
	open_time       TIME,
	close_time      TIME
);

CREATE TABLE IF NOT EXISTS way (
	id              SERIAL PRIMARY KEY,
	client_id       INT,
	chart_id        INT,
	title           VARCHAR(255),
	duration        INTERVAL,
	init_date       TIMESTAMPTZ
);

CREATE TABLE IF NOT EXISTS way_item (
	id              SERIAL PRIMARY KEY,
	way_id          INT,
	nexus           nexus_type,
	step_nomer      INT
);

CREATE TABLE IF NOT EXISTS chart_branch (
	id              SERIAL PRIMARY KEY,
	chart_id        INT,
	branch_id       INT
);

CREATE TABLE IF NOT EXISTS branch_station (
	id              SERIAL PRIMARY KEY,
	branch_id       INT,
	station_id      INT
);

CREATE TABLE IF NOT EXISTS railway (
	id              SERIAL PRIMARY KEY,
	from_id         INT,
	to_id           INT,
	duration        INTERVAL
);

CREATE TABLE IF NOT EXISTS station_transition (
	id              SERIAL PRIMARY KEY,
	station_id      INT,
	transition_id   INT
);

CREATE TABLE IF NOT EXISTS way_item_station (
	id              SERIAL PRIMARY KEY,
	way_item_id     INT,
	station_id      INT
);

CREATE TABLE IF NOT EXISTS way_item_railway (
	id              SERIAL PRIMARY KEY,
	way_item_id     INT,
	railway_id      INT
);
\end{lstlisting}

\clearpage

\begin{lstlisting}[label=lst:tbl3,caption=Запросы для создания таблиц (часть 3)]
CREATE TABLE IF NOT EXISTS way_item_transition (
	id              SERIAL PRIMARY KEY,
	way_item_id     INT,
	transition_id   INT
);
\end{lstlisting}

\section{Реализация ограничений сущностей}

В листингах~\ref{lst:cnst1}--\ref{lst:cnst3} приведены запросы для создания ограницений таблиц.
Для поддержки модульности настройка ограничений была вынесена в отдельный файл.
Преимущественно в нем настраиваются внешние ключи и проверки уникальности определенных
столбцов таблиц.

\section{Реализация функций}
В листингах~\ref{lst:add_chart_json:1}--\ref{lst:del:tr:funcs}
приведены реализации функций, где в листингах~\ref{lst:add_chart_json:1}~--~\ref{lst:add_chart_json:2}~и~\ref{lst:add_route_json:1}~--~\ref{lst:add_route_json:4}
представлены функции добавления схемы и маршрута соответственно,
которые передаются как аргументы в json формате.

В листингах~\ref{lst:get_chart_json_by_id:1}~--~\ref{lst:get_route_json_by_id:5}
приведено получение схемы и маршрута по айди в jsonb формате.

В листингах~\ref{lst:color_convertors},~\ref{lst:nav_funcs}
представлены вспомогательные функции конвертации цвета и поиска айди сущностей
по параметрам соответственно.

В листингах~\ref{lst:ck:tr:funcs:1}~--~\ref{lst:del:tr:funcs}~--~триггерные функции.
Они небходимы по вдум причинам. Для поддержания целостности базады данных. То есть,
при добавлении новой схемы, они помимо наложенных ограничений, гарантирую корректность
связывания таблиц:

\begin{itemize}
	\item[---] Проверка уникальной связи вида схема-ветка, а именно одна и та же ветка может быть дочерней по отношению только к одной схеме;
	\item[---] Проверка уникальной связи вида ветка-станция, а именно одна и та же станция может быть дочерней по отношению только к одной ветке;
	\item[---] Проверка единственной связи вида станция-переезд-станция, а именно две станции могут быть связаны только одним переездом;
	\item[---] Проверка локальности связи вида станция-переезд-станция, то есть станция, связанные переездом должны принадлежать одной и той же ветке;
	\item[---] Проверка локальности связи вида станция-переход-станция, то есть станции, связанные переходом, должны принадлежать разным веткам, которые в свою очередь принадлежат одной и той же схеме.
\end{itemize}

Также триггеры обеспечивают корректное удаление записей из таблиц. Ввиду имения сложных связывающий таблиц,
представляющие отдельные сущности, каскадное удаление не позволяет найти все взаимосвязи, поэтому требуется
вручную искать связынные записи и удалять их.

\section{Реализация ролевой модели}

В листингах \ref{lst:roles:1}--\ref{lst:roles:2} приведена реализация ролевой модели.

\begin{lstlisting}[label=lst:roles:1,caption=Реализация ролей (часть 1),firstnumber=1]
REVOKE CREATE ON SCHEMA public FROM PUBLIC;

-- UNSIGNED CLIENT --

CREATE ROLE unsigned_client WITH LOGIN PASSWORD 'qwer-tyui';

ALTER ROLE unsigned_client 
	WITH NOCREATEDB NOCREATEROLE NOREPLICATION NOINHERIT
	VALID UNTIL 'infinity';

GRANT USAGE ON SCHEMA public TO unsigned_client;

GRANT SELECT ON
	client,
	chart,
	chart_branch,
	branch,
	branch_station,
	station,
	station_transition,
	transition,
	railway
TO unsigned_client;

GRANT EXECUTE ON FUNCTION
	get_chart_json_by_id(INT),
	get_station_id_by_branch_id_station_title(INT, TEXT),
	get_station_id_by_branch_title_station_title(TEXT, TEXT),
	hex_color_to_int(TEXT),
	int_color_to_hex(INT)
TO unsigned_client;

-- SIGNED CLIENT --

CREATE ROLE signed_client WITH LOGIN PASSWORD 'asdf-ghjk';

ALTER ROLE signed_client 
	WITH NOCREATEDB NOCREATEROLE NOREPLICATION INHERIT
	VALID UNTIL 'infinity';

GRANT unsigned_client TO signed_client;
\end{lstlisting}

\clearpage

\begin{lstlisting}[label=lst:roles:2,caption=Реализация ролей (часть 2),firstnumber=43]
GRANT SELECT, INSERT, UPDATE, DELETE ON
	client,
	way,
	way_item,
	way_item_railway,
	way_item_station,
	way_item_transition
TO signed_client;

GRANT USAGE ON SEQUENCE
	client_id_seq,
	way_id_seq,
	way_item_id_seq,
	way_item_railway_id_seq,
	way_item_station_id_seq,
	way_item_transition_id_seq
TO signed_client;

GRANT EXECUTE ON FUNCTION
	add_chart_json(TEXT),
	add_route_json(INT, INT, TEXT),
	get_route_json_by_id(INT)
TO signed_client;

-- DUTY --

CREATE ROLE duty WITH LOGIN PASSWORD 'zxcv-bnml';

ALTER ROLE duty 
	WITH NOCREATEDB NOCREATEROLE NOREPLICATION INHERIT
	VALID UNTIL 'infinity';

GRANT signed_client TO duty;

GRANT UPDATE ON
	station,
	transition,
	railway
TO duty;
\end{lstlisting}

\clearpage
